\documentclass[10pt,a4paper]{amsart}
\usepackage[margin=19mm]{geometry}
\usepackage{amsmath,amssymb,mathtools}
\usepackage[scr=rsfs,cal=euler]{mathalfa}
\usepackage{xfrac}
\usepackage[unicode,hidelinks,bookmarks=false]{hyperref}
\newcommand{\ie}{i.e.\ }
\newcommand{\cf}{cf.\ }
\newcommand{\resp}{respectively\ }
\newcommand{\Subsection}{\subsection}
\providecommand{\nobreakdash}{\nobreak\hbox{-}}
\setlength{\emergencystretch}{2em}
\multlinegap0pt
\usepackage{xcolor}
\usepackage{amssymb}
\usepackage{mathtools}
\usepackage{bm}
\usepackage{enumerate}
\newtheorem{theorem}{Theorem}
\newcommand{\Cinfty}{\mathscr{C}^\infty}
\DeclareMathOperator{\Ima}{Im}
\newcommand{\dd}{\mathrm{d}}
\newcommand\restr[2]{\left.#1\right|_{#2}} \newcommand{\bfX}{\mathbf{X}} \newcommand{\bfY}{\mathbf{Y}} \newcommand{\bfZ}{\mathbf{Z}} \newcommand{\qeddiamond}{\hfill $\diamond$} \newcommand{\parder}[2]{\frac{\partial #1}{\partial #2}} \newcommand{\dparder}[2]{\dfrac{\partial #1}{\partial #2}} \newcommand{\tparder}[2]{\partial #1/\partial #2} \newcommand{\parderr}[3]{\frac{\partial^2 #1}{\partial #2\partial #3}} \newcommand{\dparderr}[3]{\dfrac{\partial^2 #1}{\partial #2\partial #3}} \newcommand{\tparderr}[3]{\partial^2 #1/\partial #2\partial #3}
\title{Hamilton--Jacobi theory for non-conservative field theories in the $k$-contact framework}
\author{Javier de Lucas, Julia Lange, Xavier Rivas, Cristina Sard\'on}
\date{Source: arXiv:2604.27670v1 (CC BY 4.0)}
\begin{document}
\maketitle

\noindent\emph{Source and licence.} Hamilton--Jacobi theory for non-conservative field theories in the $k$-contact framework. Version arXiv:2604.27670v1 (CC BY 4.0), distributed by arXiv under the Creative Commons Attribution 4.0 International licence, http://creativecommons.org/licenses/by/4.0/, as recorded in the arXiv licence metadata for this exact version and shown on the abstract page https://arxiv.org/abs/2604.27670v1. Full licence notice: \texttt{licenses/arxiv-2604.27670-CC-BY-4.0.txt}.\par\smallskip

\noindent\textit{Editorial scope.} Counted unit: the article's theorem theo:evolutionHJzindependent (source0.tex lines 1013--1019). The article's complete original proof of it is delivered with it (source0.tex lines 1021--1168, one \texttt{proof} environment of 148 lines). No mathematics is written, repaired or re-derived: the statement and the proof are the article's own lines, in the article's own notation, and no retained line is substituted or re-flowed.  Retained uncounted as the source-local context the proof needs: lines 644--650; lines 855--857; lines 913--916.  Nothing was omitted from any retained span, so the package carries no omission record.  No retained line contains a citation, so the wrapper carries no bibliography and no bibliography entry.  No retained line contains a figure environment, an image-inclusion command or drawing code, so no image is delivered and the package declares no asset dependency.  Editorial formatting only, in addition: an AMS \texttt{amsart} preamble that restates the article's own declarations and macros used by the retained lines, namely the environments \texttt{theorem} on separate counters, together with the standard packages those lines need.  The retained environments print in this document as Theorem 1 in order of appearance, whereas the article prints the same items with the numbering of its own full document.  The complete native source of the article as distributed in the arXiv e-print for this exact version is bundled unchanged as \texttt{sources/199492/source0.tex}.
\begin{equation}\label{eq:k-contact-HdDW-fields-Darboux-coordinates}
\begin{dcases}
(E_\beta)^i = \parder{h}{p_i^\beta}\,,\\
\sum_{\alpha=1}^k  (E_\alpha)^\alpha_i = -\left( \parder{h}{q^i} + \sum_{\alpha=1}^k  p_i^\alpha\parder{h}{z^\alpha} \right)\,,\\
\sum_{\alpha=1}^k  (E_\alpha)^\alpha = \sum_{j=1}^n  \sum_{\alpha=1}^k  p_j^\alpha\parder{h}{p_j^\alpha}\,.
\end{dcases}\qquad i=1,\dotsc,n\,,\qquad \beta=1,\dotsc,k\,.
\end{equation}

\begin{equation}\label{eq:ProjZgamma}
({\bf X}^\gamma)_\alpha(q)=\sum_{i=1}^n\frac{\partial h}{\partial p^\alpha_i}\circ\gamma(q)\frac{\partial}{\partial q^i}\,,\qquad \alpha=1,\dotsc,k\,,\qquad q\in Q\,.
\end{equation}

\begin{equation}\label{eq:closedgamma2}
\frac{\partial \gamma_i^\mu}{\partial q^j} =
\frac{\partial \gamma_j^\mu}{\partial q^i}\,,\qquad \mu=1,\dotsc,k\,,\qquad 1\leq i<j\leq n\,,
\end{equation}

\begin{theorem}\label{theo:evolutionHJzindependent}
Let $\mathbf{E}_h$ be an ${\bm \eta}_{Q,k}$-evolution $k$-vector field associated with $h\in \Cinfty(\mathcal{J}_{Q,k})$, let $\gamma$ be a holonomic section, and assume $\mathbf{E}_h^\gamma$ to be integrable. Then, the following statements are equivalent:
\begin{enumerate}
\item Every integral section $\sigma$ of $\mathbf{E}_h^\gamma$ lifts to a solution $\gamma\circ\sigma$ of the $z$-independent ${\bm \eta}_{Q,k}$-evolution HdDW equations,
\item $\dd(h\circ\gamma)=0$.
\end{enumerate}
\end{theorem}

\begin{proof}
Suppose that $(\gamma \circ \sigma)(t) = (\sigma^i(t), \gamma_i^\alpha(\sigma(t)), \gamma^\alpha(\sigma(t)))$ is a solution to the $z$-independent HdDW equations for ${\bm \eta}_{Q,k}$-evolution  $k$-vector fields. Then,
\[
\begin{dcases}
\restr{\dfrac{\partial \sigma^i}{\partial t^\alpha}}{t} = \restr{\dfrac{\partial h}{\partial p_i^\alpha}}{\gamma\circ \sigma(t)}\,,\\[2mm]
\sum_{\alpha=1}^k \restr{\dfrac{\partial(\gamma_i^\alpha\circ \sigma)}{\partial t^\alpha}}{t}
=
-\left( \restr{\dfrac{\partial h}{\partial q^i}}{\gamma\circ \sigma(t)} + \sum_{\alpha=1}^k\left(\gamma_i^\alpha \circ \sigma(t) \right) \restr{\dfrac{\partial h}{\partial z^\alpha}}{\gamma\circ \sigma(t)}\right)\,,  \\[2mm]
\sum_{\alpha=1}^k \restr{\dfrac{\partial(\gamma^\alpha\circ \sigma)}{\partial t^\alpha}}{t}
=
\sum_{\alpha=1}^k\sum_{j=1}^n (\gamma^\alpha_j\circ \sigma)(t) \restr{\dfrac{\partial h}{\partial p_j^\alpha}}{\gamma\circ \sigma(t)},
\end{dcases}
\]
for $\alpha=1,\dotsc,k,\quad i=1,\dotsc,n\,.$ On the other hand,
\[
\dd(h \circ \gamma)
=
\sum_{i=1}^n\left( \frac{\partial h}{\partial q^i} \circ \gamma + \sum_{\alpha=1}^k\sum_{j=1}^n\left( \frac{\partial h}{\partial p_j^\alpha}\circ \gamma \right) \frac{\partial \gamma_j^\alpha}{\partial q^i} + \sum_{\alpha=1}^k\left( \frac{\partial h}{\partial z^\alpha}\circ \gamma \right) \frac{\partial \gamma^\alpha}{\partial q^i} \right) \dd q^i.
\]
Since $\gamma$ is holonomic, it follows that \eqref{eq:closedgamma2} holds.
Using \eqref{eq:closedgamma2} and the $z$-independent $\bm\eta_{Q,k}$-evolution HdDW equations, we have
\begin{align*}
\restr{\dd(h\circ \gamma)}{\sigma(t)}
&=
\sum_{i=1}^n
\Bigg(
  \restr{\frac{\partial h}{\partial q^i}}{\gamma(\sigma(t))}
  + \sum_{\alpha=1}^k \sum_{j=1}^n
    \restr{\frac{\partial h}{\partial p^\alpha_j}}{\gamma(\sigma(t))}
    \restr{\frac{\partial \gamma^\alpha_j}{\partial q^i}}{\sigma(t)}\\[1ex]&
  + \sum_{\alpha=1}^k
    \restr{\frac{\partial h}{\partial z^\alpha}}{\gamma(\sigma(t))}
    \restr{\frac{\partial \gamma^\alpha}{\partial q^i}}{\sigma(t)}
\Bigg)
\, \dd q^i(\sigma(t))
\\[1ex]
&=
\sum_{i=1}^n
\Bigg(
  - \sum_{\alpha=1}^k
    \left(\gamma_i^\alpha\circ \sigma\right)(t)
    \restr{\frac{\partial h}{\partial z^\alpha}}{\gamma(\sigma(t))}
  - \sum_{\alpha=1}^k
    \restr{\frac{\partial(\gamma_i^\alpha\circ\sigma)}{\partial t^\alpha}}{t}
\\
&\qquad\qquad
  + \sum_{\alpha=1}^k \sum_{j=1}^n
    \restr{\frac{\partial \sigma^j}{\partial t^\alpha}}{t}
    \restr{\frac{\partial \gamma^\alpha_i}{\partial q^j}}{\sigma(t)}
  + \sum_{\alpha=1}^k
    \restr{\frac{\partial h}{\partial z^\alpha}}{\gamma(\sigma(t))}
    \restr{\frac{\partial \gamma^\alpha}{\partial q^i}}{\sigma(t)}
\Bigg)
\, \dd q^i(\sigma(t))
\\[1ex]
&=
\sum_{i=1}^n
\sum_{\alpha=1}^k
\left(
  - \gamma_i^\alpha(\sigma(t))
  + \restr{\frac{\partial \gamma^\alpha}{\partial q^i}}{\sigma(t)}
\right)
\restr{\frac{\partial h}{\partial z^\alpha}}{\gamma(\sigma(t))}
\, \restr{\dd q^i}{\sigma(t)} = 0\,,
\end{align*}
where we used in the last line the fact that $\dfrac{\partial \gamma^\alpha}{\partial q^i} = \gamma^\alpha_i$ for all possible indices. Hence, $\dd(h\circ \gamma)(q) = 0$ for every $q\in \Ima \sigma$. Since $\mathbf{E}_h^\gamma$ is integrable by assumption, every $q\in Q$ induces an integral section $\sigma$ such that $\sigma(0) = q$. Therefore, $\dd(h\circ\gamma)(q) = 0$ for every $q\in Q$.

Let us prove the converse. Assume that $\dd (h \circ \gamma) = 0$, and let $\sigma$ be an integral section of $\mathbf E_h^\gamma$. Let us prove that $\gamma \circ \sigma$ is a solution to the $z$-independent ${\bm \eta}_{Q,k}$-evolution HdDW equations. From $\dd(h\circ \gamma)=0$, we have
\begin{equation}
\label{eq:d(hgamma)=0-evol}
0
=
\sum_{i=1}^n
\Bigg(
  \frac{\partial h}{\partial q^i} \circ \gamma
  + \sum_{\alpha=1}^k \sum_{j=1}^n
    \left(\frac{\partial h}{\partial p^\alpha_j} \circ \gamma \right)
    \frac{\partial \gamma^\alpha_j}{\partial q^i}
  + \sum_{\alpha=1}^k
    \left(\frac{\partial h}{\partial z^\alpha} \circ \gamma \right)
    \frac{\partial \gamma^\alpha}{\partial q^i}
\Bigg)
\, \dd q^i.
\end{equation}
Since $\mathbf E_h$ is an evolution ${\bm \eta}_{Q,k}$-Hamiltonian $k$-vector field, equations \eqref{eq:ProjZgamma} and \eqref{eq:k-contact-HdDW-fields-Darboux-coordinates} yield
\[
(\mathbf E_h^\gamma)_\alpha
=
\sum_{i=1}^n
\left(
  \frac{\partial h}{\partial p_i^\alpha}
  \circ \gamma
\right)
\frac{\partial}{\partial q^i}\,,
\qquad \alpha=1,\dotsc,k\,.
\]
Since $\sigma$ is an integral section of ${\bf E}_h^\gamma$, it follows that
\begin{equation}
\label{eq:sigmaHDW-evol}
\restr{\frac{\partial \sigma^i}{\partial t^\alpha}}{t} = \frac{\partial h}{\partial p_i^\alpha}\circ \gamma \circ \sigma(t)\,,\qquad i=1,\dotsc,n\,,\qquad \alpha=1,\dotsc,k\,.
\end{equation}
From \eqref{eq:closedgamma2}, \eqref{eq:d(hgamma)=0-evol}, and \eqref{eq:sigmaHDW-evol}, one has
\begin{align}
\sum_{\alpha=1}^k
\restr{\frac{\partial(\gamma_i^\alpha \circ \sigma)}{\partial t^\alpha}}{t}
&=
\sum_{\alpha=1}^k \sum_{j=1}^n
\restr{\frac{\partial \gamma_i^\alpha}{\partial q^j}}{\sigma(t)}
\restr{\frac{\partial \sigma^j}{\partial t^\alpha}}{t}
\\
&=
\sum_{\alpha=1}^k \sum_{j=1}^n
\restr{\frac{\partial \gamma_i^\alpha}{\partial q^j}}{\sigma(t)}
\restr{\frac{\partial h}{\partial p_j^\alpha}}{\gamma(\sigma(t))}
=
\sum_{\alpha=1}^k \sum_{j=1}^n
\restr{\frac{\partial \gamma_j^\alpha}{\partial q^i}}{\sigma(t)}
\restr{\frac{\partial h}{\partial p_j^\alpha}}{\gamma(\sigma(t))}
\\
&=
- \restr{\frac{\partial h}{\partial q^i}}{\gamma(\sigma(t))}
-
\sum_{\alpha=1}^k
\restr{\frac{\partial \gamma^\alpha}{\partial q^i}}{\sigma(t)}
\restr{\frac{\partial h}{\partial z^\alpha}}{\gamma(\sigma(t))}
\\
&=
- \restr{\frac{\partial h}{\partial q^i}}{\gamma(\sigma(t))}
-
\sum_{\alpha=1}^k
\gamma_i^\alpha (\sigma(t))\,
\restr{\frac{\partial h}{\partial z^\alpha}}{\gamma(\sigma(t))}.
\end{align}
Moreover,
\[
\sum_{\alpha=1}^k
\restr{\frac{\partial (\gamma^\alpha \circ \sigma)}{\partial t^\alpha}}{t}
=
\sum_{\alpha=1}^k \sum_{i=1}^n
\restr{\frac{\partial \gamma^\alpha}{\partial q^i}}{\sigma(t)}
\restr{\frac{\partial \sigma^i}{\partial t^\alpha}}{t}
=
\sum_{\alpha=1}^k \sum_{i=1}^n
\gamma^\alpha_i(\sigma(t))\,
\restr{\frac{\partial h}{\partial p_i^\alpha}}{\gamma\circ\sigma(t)}.
\]
This concludes the proof.
\end{proof}

\end{document}
